\documentclass[10pt,a4paper]{amsart}
\usepackage[margin=19mm]{geometry}
\usepackage{amsmath,amssymb,mathtools}
\usepackage[scr=rsfs,cal=euler]{mathalfa}
\usepackage{xfrac}
\usepackage[unicode,hidelinks,bookmarks=false]{hyperref}
\newcommand{\ie}{i.e.\ }
\newcommand{\cf}{cf.\ }
\newcommand{\resp}{respectively\ }
\newcommand{\Subsection}{\subsection}
\providecommand{\nobreakdash}{\nobreak\hbox{-}}
\setlength{\emergencystretch}{2em}
\multlinegap0pt
\usepackage{bm}
\usepackage{bbm}
\usepackage{dsfont}
\usepackage{xspace}
\usepackage{cancel}
\usepackage{mathtools}
\usepackage{amsthm}
\makeatletter
\@ifundefined{theorem}{\newtheorem{theorem}{Theorem}}{}
\makeatother
\newcommand{\BI}{\mathcal{BI}}
\newcommand{\Baire}{\omega^{\omega}}
\newcommand{\FF}{\mathcal{F}in\times\mathcal{F}in}
\newcommand{\aaa}{\mathfrak{a}}
\newcommand{\bbb}{\mathfrak{b}}
\title[On antichain numbers and the splitting ideal]{On antichain numbers and the splitting ideal}
\author{Aleksander Cie\'slak}
\date{Source: arXiv:2605.19044v2 (CC BY 4.0)}
\begin{document}
\maketitle

\noindent\emph{Source and licence.} On antichain numbers and the splitting ideal. Version arXiv:2605.19044v2 (CC BY 4.0), distributed under the Creative Commons Attribution 4.0 International licence, as recorded in the arXiv licence metadata for this exact version and shown on the abstract page https://arxiv.org/abs/2605.19044v2. Full rights notice: licenses/arxiv-2605.19044-CC-BY-4.0.txt.\par\smallskip

\noindent\textit{Editorial scope.} Counted unit: the Theorem whose source label is \texttt{bbbLeqAAABi} in the article (source0.tex lines 767--769, source label \texttt{bbbLeqAAABi}), delivered together with the article's own complete original proof of it (source0.tex lines 770--809, one \texttt{proof} environment of 40 lines). No mathematics is written, repaired or re-derived: statement and proof are the article's own text line for line. Required context retained uncounted: none. Every definition, display and label of those lines is retained unchanged. Omitted from this package and not redistributed: 1--766, 810--1071. Nothing inside a retained excerpt is omitted or altered: every retained body is delivered exactly as the article's lines are written, so no omission receipt of any kind accompanies this package. Citations: the retained excerpts carry no citation command at all, so no bibliography entry of the article is transcribed and no reference is redistributed. Reference adaptation: none was needed and none was made; every reference inside the delivered unit resolves inside it, so no unresolved reference or citation is produced. Printed slips kept literal: no printed word, symbol, hypothesis or number of the counted statement or of its proof was altered. Layout and class: the article's own class is \texttt{amsart} with packages including amsmath, amsfonts, amssymb, amsthm, mathrsfs, yfonts, enumitem, tikz, tikz-cd, amsaddr, comment; the wrapper is the lane's pinned \texttt{amsart} editorial wrapper at 10pt on A4, which restates the theorem-like environments the retained text needs and adds the article's own macro definitions that the retained text needs (\texttt{\textbackslash BI}, \texttt{\textbackslash Baire}, \texttt{\textbackslash FF}, \texttt{\textbackslash aaa}, \texttt{\textbackslash bbb}), with extra wrapper packages (bm, bbm, dsfont, xspace, cancel, mathtools). The delivered numbering is the unit's own. Assets: no retained span contains a figure environment, a graphics-inclusion command, a PDF-page inclusion, a table or a TikZ picture; no image file is copied into this package, the package declares no asset dependency and no image file is redistributed with it. Licence: version arXiv:2605.19044v2 of this article is distributed under the Creative Commons Attribution 4.0 International licence, as recorded in the arXiv licence metadata for this exact version and shown on the abstract page https://arxiv.org/abs/2605.19044v2. A full rights notice is delivered at licenses/arxiv-2605.19044-CC-BY-4.0.txt. Characters: every retained excerpt is pure ASCII, so the delivered engine, XeLaTeX with its default Latin Modern text font, reports no missing character.
\begin{theorem}\label{bbbLeqAAABi}
    $\bbb\leq\aaa(\BI)$
\end{theorem}

\begin{proof}
    Assume that $\kappa<\bbb$ and that a $\BI$-antichain $\{A_{\alpha}:\alpha<\kappa\}\subseteq\BI^{+}$ is given.\\
    We will say that $\alpha<\kappa$ is of type 1, if there is an $i_{\alpha}\in\omega$ such that the set $\{(j,k):(i_{\alpha},j,k)\in A_{\alpha}\}$ is $\FF$-positive. In such case the set $\{(j,k):(i_{\alpha},j,k)\in A_{\alpha}\}$ contains a set of the form $F_{\alpha}=\bigcup_{n}\{j^{\alpha}_{n}\}\times F^{\alpha}_{n}$ where $\{j^{\alpha}_{n}:n\in\omega\}\in[\omega]^{\omega}$ and $F^{\alpha}_{n}\in[\omega]^{\omega}$ for every $n\in\omega$.\\
    We will say that $\alpha<\kappa$ is of type 2, if it is not of type 1. Note that if $\alpha$ is of type 2, then the set $I_{\alpha}=\{i\in\omega:|(\{i\}\times\omega^{2})\cap A_{\alpha}|=\omega\}$ is infinite.\\
    
    Consider now the following two cases:\\
    CASE 1: the set $W=\{\alpha<\kappa: \alpha$ is of type 1$\}$ is uncountable.\\
    Without loss of generality we may assume that $\omega\subseteq W$ and $i_{\alpha}$'s for $\alpha\in W$ are equal to the same $i^{*}\in\omega$. For every $\alpha<\kappa$ and $l\in\omega$ with $l\neq\alpha$ there are $K^{\alpha}(l)\in\omega$ and $f^{\alpha}_{l}\in\Baire$ such that:
    \begin{center}
    $\{(j,k)\in\omega^{2}:(i^{*},j,k)\in A_{l}\cap A_{\alpha}\}\subseteq\{(j,k)\in\omega^{2}:j<K^{\alpha}(l)$ or $k\leq f^{\alpha}_{l}(j)\}$
    \end{center}
    As $\kappa<\bbb$ there is $f\in\Baire$ that dominates all $f^{\alpha}_{l}$'s. Then, for every $\alpha\in\kappa\setminus\omega$ we define $g_{\alpha}\in\Baire$ in a way that for all $j\geq g_{\alpha}(l)$ we have $f^{\alpha}_{l}(j)<f(j)$. Again, as $\kappa<\bbb$, we can find $g\in\Baire$ that dominates all functions $g_{\alpha}$ and $K^{\alpha}$.\\
    Inductively we construct sequences $\{j_{l}:l\in\omega\}$ and $\{C_{l}:l\in\omega\}$ such that for every $l\in\omega$ we have:
    \begin{itemize}
        \item[--] $j_{l+1}>j_{l}>g(l)$,
        \item[--] $C_{l}=F_{l}\cap (\{j_{l}\}\times(\omega\setminus f(j_{l}))$ is infinite
    \end{itemize}
    When the induction is finished define $C=\bigcup_{l}\{j_{l}\}\times C_{l}$ and $A=\{i^{*}\}\times C$. Note that the set $A$ is $\BI$-positive because $C$ is $\FF$-positive. We claim that for every $\alpha<\kappa$ we have $A\cap A_{\alpha}\in\BI$ which will finish the first case. Fix $\alpha<\kappa$. If $\{(j,k)\in\omega^{2}:(i^{*},j,k)\in A_{\alpha}\}\in\FF$, there is nothing to do. Assume then, that $\{(j,k)\in\omega^{2}:(i^{*},j,k)\in A_{\alpha}\}$ is $\FF$-positive. Let $L\in\omega$ be large enough so that for all $l>L$ we have $g(l)\geq\max\{g_{\alpha}(l),K^{\alpha}(l)\}$. But then
    \begin{center}
        $\{(j,k)\in\omega^{2}:(i^{*},j,k)\in A\cap A_{\alpha}\}\subseteq\bigcup_{l<L}C_{l}$
    \end{center}
    This is because if $l>L$, then $j_{l}>\max\{g_{\alpha}(l),K^{\alpha}(l)\}$ and we have that
    \begin{center}
    $C_{l}\cap\{(j,k)\in\omega^{2}:(i^{*},j,k)\in A_{\alpha}\}\subseteq C_{l}\cap \{(j,k):j<K^{\alpha}(l)$ or $k<f^{\alpha}_{l}(j)\}$
    \end{center}
    Notice that the latter is empty because $C_{l}\subseteq\{j_{l}\}\times(\omega\setminus f(j_{l}))$ and $f(j_{l})>f^{\alpha}_{l}(j_{l})$ for all $j_{l}>K^{\alpha}(l)$. As $\bigcup_{l<L}C_{l}\in\FF$, the first case is finished.\\
    
    CASE 2: $W$ is countable.\\
    Let $\{A'_{n}:n\in\omega\}$ enumerate all of the sets $\{A_{\alpha}:\alpha\in W\}$. Then, by removing the set $\{A'_{n}:n\in\omega\}$ from the original enumeration, we may assume that the intersection $\{A'_{n}:n\in\omega\}\cap\{A_{\alpha}:\alpha<\kappa\}$ is empty.\\
    For every $\alpha\in\kappa\setminus\omega$ define $g_{\alpha}\in\Baire$ such that for every $l\in\omega$ for all $i\geq g_{\alpha}(l)$ the set $\{(j,k):(i,j,k)\in A_{l}\cap A_{\alpha}\}$ is finite. Then, for all $\alpha<\kappa$ and $l\in\omega$ with $\alpha\neq l$ define $f^{\alpha}_{l}\in\Baire$ such that $\{(j,k):(i,j,k)\in A_{l}\cap A_{\alpha}\}\subseteq f^{\alpha}_{l}(i)$ for all $i>g_{\alpha}(l)$.
    As $\kappa<\bbb$ we can find $f\in\Baire$ that dominates all such $f^{\alpha}_{l}$'s.
    For every $\alpha\in\kappa\setminus\omega$ define then $h_{\alpha}\in\Baire$ such that for every $i>h_{\alpha}(l)$ we have $f(i)>f^{\alpha}_{l}(i)$. As $\kappa<\bbb$ there is $g\in\Baire$ that dominates all $g_{\alpha}$'s and $h_{\alpha}$'s.\\
    Inductively we construct an increasing sequence $\{i_{l}:l\in\omega\}\subseteq\omega$ and $\{C_{l}:l\in\omega\}$ such that for every $l\in\omega$ we have:
    \begin{itemize}
        \item[--] $i_{l}\in I_{l}\setminus i_{l-1}$,
        \item[--] $C_{l}=\{(j,k):(i_{l},j,k)\in A_{l}\setminus(\bigcup_{l'\leq l}A'_{l'}\cup f(i_{l}))\}$
        \item[--] $i_{l}>g(l)$
    \end{itemize}
    When induction is over define $A=\bigcup_{l}\{i_{l}\}\times C_{l}$. Note that $A$ is $\BI$-positive as each $C_{l}$ is infinite and $\{i_{l}:l\in\omega\}$ is increasing. By a similar argument as in the first case, we have that $A\cap A_{\alpha}\in\BI$ for every $\alpha\in\kappa$ and $A\cap A'_{n}\in\BI$ for every $n\in\omega$. This finishes the proof.
\end{proof}

\end{document}
